\section{Which difficulties are resolved, and which data are still missing}
\label{sec:applications}

The new results make two kinds of progress. They turn certain informal
resource arguments into actual PR programs, and they turn particular
geometric certificates into complete point lists. They do not supply
every certificate that an existing application would need. This section
states the resulting interfaces precisely.

\subsection{Finite exhaustive witnesses as a general interface}

The auxiliary-polynomial construction is an instance of a useful
general rule. It applies to finite chart families, bounded algebraic
certificates, bounded collections of exact matrices, and finite
geometric decompositions, provided each candidate has a certified
complete finite output.

\begin{theorem}[Elimination of a bounded auxiliary choice]
\label{thm:exhaustive}
Let \(P(a,p)\), \(\Good(a,c)\), and \(R(a,c,p)\) be fixed PR
predicates on natural codes. Let \(B(a)\) and the finite list
\(L(a,c)\) be fixed PR programs. Suppose \(\ipra\) proves, for every
good candidate,
\[
 \Good(a,c)\ \longrightarrow\
 \bigl(p\in L(a,c)\ \longleftrightarrow\ R(a,c,p)\bigr).
\]
Suppose also that it proves the pointwise coverage assertion
\begin{equation}\label{eq:bounded-witness-cover}
 P(a,p)\ \longrightarrow\
 \exists c\le B(a)\,
       \bigl(\Good(a,c)\wedge R(a,c,p)\bigr).
\end{equation}
Then a fixed PR program produces a duplicate-free complete list of
all \(p\) satisfying \(P(a,p)\), and \(\ipra\) proves its completeness.
No chosen function sending \(p\) to a candidate \(c\) is required.
\end{theorem}

\begin{proof}
Scan \(c=0,\ldots,B(a)\). For good candidates concatenate \(L(a,c)\),
retain precisely entries satisfying \(P(a,\cdot)\), and remove
duplicates by a bounded scan through the accumulated list.
All loops have PR bounds: the outer counter is explicit and each
inner list has a PR length. The scan invariant says, for a fixed
parameter \(p\), that membership in the accumulated list is
equivalent to
\[
 P(a,p)\wedge\exists c<i\,
       (\Good(a,c)\wedge p\in L(a,c)).
\]
This is a bounded formula; its quantifier-free compilation is the
induction formula. At the end, every retained entry satisfies \(P\)
by the filter. Conversely, given \(P(a,p)\), eliminate the existential
in \eqref{eq:bounded-witness-cover}. Its witness is an enumerated good
candidate, and the candidate-completeness theorem puts \(p\) into
the corresponding list. The invariant places it in the final list.
Universal introduction over \(p\) follows after the finite induction.
\end{proof}

If \(P\) is not decidable, omit the filter. The same argument yields
a finite list containing every target point, hence a height or code
bound, but not necessarily an exact list of target points. This is
often the correct output contract when the remaining property is
not an arithmetic Boolean test.

Theorem~\ref{thm:exhaustive} does not require an explicit bound on
the \emph{number} of successful candidates smaller than \(B(a)+1\).
It relies instead on a finite bound on their possible codes. This is
why a coefficient bound on auxiliary equations can do more than a
cardinality bound on exceptional subspaces.

\subsection{A precise obstruction to converting cardinality into a list}

The distinction above is not only a bookkeeping concern. There is an
elementary computability counterexample, already anticipated by the
repository's height/list warnings~\cite{repo-faltings,repo-framework}.

\begin{proposition}[Uniform singleton bounds do not locate their member]
\label{prop:singleton}
There is one fixed PR predicate \(P(e,t)\) such that \(\ipra\) proves
that every duplicate-free list satisfying \(P(e,\cdot)\) has length
at most one, uniformly in \(e\), but there is no total computable
procedure producing a complete finite list for every \(e\). Nor is
there a total computable bound \(H(e)\) such that
\(P(e,t)\to t\le H(e)\) for all \(e,t\).
\end{proposition}

\begin{proof}
Fix an ordinary effective enumeration of machines and let \(P(e,t)\)
say that machine \(e\) first halts at step \(t\). A bounded machine
simulation is PR, and ``first'' is a bounded test. If \(s<t\) and
\(P(e,s)\), then the negative earlier-halting condition in \(P(e,t)\)
fails. Decidable comparison of \(s,t\) therefore proves
\(P(e,s)\wedge P(e,t)\to s=t\) in \(\ipra\). Applying that fact
to two distinct positions proves the list bound.

A total computable complete-list procedure would decide whether the
machine halts by testing whether its output list is empty. A total
computable bound \(H(e)\) would decide the same question by running
the machine through that many steps and checking the finite trace.
Either conclusion contradicts the elementary undecidability of the
halting problem. The simulator is bounded; no total universal
primitive-recursive evaluator is being postulated.
\end{proof}

The example can be placed directly inside the auxiliary-polynomial
interface. Take the fixed curve \(F(X,Y)=Y\) and the target points
\((t,0)\) with \(P(e,t)\). Every target point admits the linear
auxiliary equation \(G_t(X,Y)=X-t\), and \(F,G_t\) have a nonzero
resultant. Every finite target list is covered by one such auxiliary
equation (use \(X\) for an empty list). Degrees are uniformly one,
and target cardinality is uniformly at most one. Nevertheless a
computable bound on the necessary coefficient heights would contradict
Proposition~\ref{prop:singleton}, because the bounded-envelope theorem
would then produce the forbidden complete-list algorithm.
Indeed, writing a linear auxiliary as \(aX+bY+c\), its good
intersection with \(Y=0\) forces \(a\ne0\) and hence
\(t=-c/a\). If the reduced coefficient heights are at most \(H\),
then \(|a|\ge1/H\) and \(|c|\le H\), giving \(t\le H^2\).
Thus the missing coefficient bound would directly locate the halting
time. This example shows exactly what a degree-only cover omits.

\subsection{The Siegel and Faltings routes}

The inspected Siegel report separates a finite-tuple approximation
cover from a geometric bound on the number of curve points in each
covering subspace~\cite{repo-siegel}. Earlier framework packets replace
some abstract divisor arguments by explicit norm calculations.
The plane-elimination theorem broadens that geometric part: if a
subspace's annihilating function is represented by an auxiliary plane
polynomial \(G\), and \(\Good(F,G)\) holds, then a finite list in
that intersection has at most \(\deg F\deg G\) points. Neither
smoothness nor a special weight shape is needed. A rational function
presentation must first clear denominators and separately account for
points where those denominators vanish; that conversion is an
additional contract, not implicit in the theorem.

For an actually supplied finite global cover by such \(G\)'s, one can
now compute all the intersection lists and obtain a complete target
list after a decidable filter. More generally, a uniform degree and
coefficient envelope for pointwise auxiliaries suffices. This gives
a concrete alternative target for an arithmetic-geometric proof:
produce the finite envelope, rather than a distinguished cover.

What is not supplied is a proof that every curve in the general Siegel
or Faltings classes admits the required envelope with PR bounds.
The existing tuple-local covering proposal, even if proved, need not
imply that stronger assertion. Likewise, nonvanishing of a particular
resultant is a check on one exact auxiliary; it is not a construction
of an auxiliary for every target point. General positive height
bounds and complete rational-point algorithms for those classes
remain outside this report. No unprovability or independence claim
about those general arithmetic sentences follows from this limitation.

\subsection{Mordell--Weil descent}

The inspected Mordell--Weil packet already gives explicit contracting
height descent, coefficient extraction, and bounded selection from
\emph{supplied} quotient representatives~\cite{repo-mw}. We do not
present these again as new results. The reset-counter theorem can
organize a refinement process that alternates finitely many ranked
tasks, once its local PR transition and rank contracts are available.
The plane-intersection procedure can solve zero-dimensional rational
subproblems arising in a particular presentation. Neither construction
produces the missing arbitrary rational quotient representatives or
proves a uniform generation theorem. Those are separate mathematical
inputs, not a loop scheduling issue.

\subsection{Analytic names, Puiseux data, and exact coefficient tests}

The new arithmetic programs use exact rational coefficient arrays.
This restriction matters. The predicate \(\Good(F,G)\) compares
finitely many exact rational coefficients with zero. If coefficients
were supplied only by arbitrary Cauchy names, a negative result from
a finite precision test would not certify exact zero. For example,
\(F=Y^2-X^2-a\) and \(G=Y-X\) have elimination behavior depending
on whether \(a=0\). Treating an arbitrary name for \(a\) as an
exact rational code would import a zero test that the input does not
provide.

The Puiseux obstruction report makes the stronger point that exact
singular coefficient data can encode zero detection even when each
individual represented value happens to be rational~\cite{repo-puiseux}.
The current framework respects that obstruction. It supports exact
algebraic codes, approximate observations, or local constructions with
positive separation certificates. It does not infer exact singular
Puiseux coefficients from arbitrary Cauchy data, and it does not add
root multiplicity matching or a globally continuous branch label.

The demand compiler does resolve a different analytic difficulty:
iterating a fixed operation while its required input accuracy changes
can be performed with a first-order PR program. The correctness proof
fixes the final precision and uses the finite horizon invariant.
Whether the local analytic estimate is valid remains a mathematical
lemma to prove, just as before.

\subsection{Prime-distribution constants and all-future claims}

The prime-distribution packet gives exact arithmetic targets for
Siegel--Walfisz and Bombieri--Vinogradov with existential constants
\emph{before} the universal range of \(X\), and isolates the missing
effective uniform information~\cite{repo-primes}. The present
framework makes finite sums, rational approximation, and their
precision schedules easier to compile. It does not move a finite
search success at each tested \(X\) in front of the universal
quantifier. Nor does the reset theorem produce a stopping rank for
searching for a constant whose correctness is an unbounded
all-future property.

The same distinction applies to metastability. A program that returns
a stable finite window for each requested control function does not
thereby return a final convergence modulus. These are different output
contracts. Our new resource theorem proves termination only after
the finite rank condition itself has an internal proof; it is not a
general mechanism for converting an unbounded mathematical promise
into a computable certificate.

\subsection{A practical dependency table}

\begin{longtable}{>{\raggedright\arraybackslash}p{0.24\textwidth}>{\raggedright\arraybackslash}p{0.34\textwidth}>{\raggedright\arraybackslash}p{0.34\textwidth}}
\toprule
Difficulty & Available resolution here & Remaining input\\
\midrule\endhead
Changing precision inside recursion
& Explicit \code{PR 2} compiler; checked semantic laws and partial internal equations
& Finish the stated arithmetic coherence proofs in Lean.\\
Lower counters reset during descent
& Fixed-dimensional lexicographic fuel, including an automatic trajectory envelope
& Internal local invariant and strict-rank proofs for the application.\\
Curve fibre count lacks monic or weight form
& Bounded shear and Sylvester fibre multiplicity, giving the degree-product bound
& A verified nonzero resultant for the actual auxiliary equation.\\
Finite auxiliaries exist pointwise but no cover was selected
& Exhaust all auxiliaries within a supplied degree and coefficient envelope
& Proof of a uniform coefficient envelope.\\
Only a tuple-local cover or a cardinality bound is known
& Correct finite-list cardinality consequence
& Positive location data for a complete list or height bound.\\
Arbitrary Cauchy coefficient equality
& Approximate or apartness-guarded interfaces remain usable
& Exact input codes or a specific separation witness.\\
General MW, Siegel, Faltings, and prime estimates
& Explicit finite components and a narrower list of obligations
& The substantive general arithmetic or analytic theorem.\\
\bottomrule
\end{longtable}
